\clearpage
\fancyhead[L]{\small BURCKHARDT. TABLES DE LA LUNE.}
\SourceLeaf{558}
\begin{center}\small Göttingische gelehrte Anzeigen. Issue 39. Pages 388..392. 8 March 1813.\\[3pt]\rule{.72\linewidth}{.3pt}\end{center}

\emph{Tables astronomiques publiées par le bureau des longitudes de France. Tables de la lune, par} M. \textsc{Burckhardt}, \emph{membre de l'institut impérial, du bureau des longitudes de France et de plusieurs autres sociétés savantes. Décembre} 1812. Paris. From Madame Courcier. 88 pages in quarto.

\textsc{Bürg}'s tables have made it possible to calculate the Moon's positions with such accuracy that it seems a daring undertaking to try to surpass these famous tables. Yet bringing the lunar tables as close to perfection as possible is so important in many respects that we owe the greatest thanks to an astronomer as skilful as the author of the present tables for his efforts to improve them still further. Work of this kind deserves all the more credit because it takes \emph{years} of sacrifice instead of a few months and, by the nature of the subject, can no longer be rewarded with \emph{brilliant} successes.

The author had two aims in undertaking this work. His first aim was to give the tables a somewhat different, more convenient, and simpler \emph{form}. To achieve it, he went beyond recasting the elements of \textsc{Bürg}'s lunar tables. He carried out a fresh analysis of more than four thousand observations and based his own tables on it, so these tables can be regarded as truly the sole property of Mr. \textsc{Burckhardt}. He could rightly hope that in this way the new tables would also have some advantage in \emph{accuracy}. This is indeed confirmed by the comparison of three hundred observations that the French Bureau of Longitudes had made with both \textsc{Burckhardt}'s and \textsc{Bürg}'s tables. For 167 longitudes observed at Greenwich and at the Imperial Observatory in Paris, the sum of the squares of the errors was $= 4602''$ with \textsc{Burckhardt}'s tables and $= 7083''$ with \textsc{Bürg}'s. For 137 other longitudes observed in Paris at the Imperial Observatory and at the Military School, the sum of the squares of the errors was $4182''$ with \textsc{Burckhardt}'s tables and $6439''$ with \textsc{Bürg}'s. The latitudes, too, according to the author's assurance, agreed better with his own tables than with \textsc{Bürg}'s. The magnitude of the deviation is not given here, however. We would have liked to see the results of all these comparisons printed individually. This would have brought home more vividly the high standard now attained by the lunar tables and made it much easier in future to see whether adding one or another new equation improves them further. In general we would have liked a little more detail in the introduction to these tables. We are by no means in favour of the lengthy instructions in which some table makers repeat familiar things until the reader is tired of them. Even so, it seems to us a reasonable wish to have the distinctive features of new tables stated briefly, fully, and explicitly, though it is easy enough to discover them by analysing the tables themselves.

The chief difference in form between these tables and the form that \textsc{Mason}, \textsc{Triesnecker}, and \textsc{Bürg} had retained since \textsc{Tobias Meyer} is the use of mean solar longitudes in place of true solar longitudes as their basis. Here the longitude of the node and the perigee are also no longer first corrected by an equation that depends on the Sun's anomaly. The evection is treated separately from the other small equations. The number of small equations whose arguments \SourceLeaf{559}can all be taken readily and independently of one another directly from the table is now 32, including the nutation and two perturbation equations due to Venus and Jupiter. These are used to correct the mean argument of the evection. The sum of those 32 equations and the evection corrects the anomaly. The sum of the 32 equations, the evection, and the equation of the centre corrects the argument of the variation. The variation, the preceding 34 equations, and the mean longitude together give the true longitude in the orbit, which is then reduced to the ecliptic. The secular equation must also be taken into account in the arguments of the evection, the equation of the centre, and the variation, and in the mean longitude. With the latter, the small long-period inequality, known to have been determined empirically, is combined in the epochs for the 19th century. For other centuries the two must be calculated separately from a special table added at the end. As for this small equation, \textsc{Laplace} now considers it most likely that it is proportional to the cosine of the sum of twice the longitude of the Moon's node and the longitude of the perigee. The period is accordingly 175 years, the coefficient is set $= 12''5$, and the equation originates in an assumed inequality between the Earth's northern and southern hemispheres. It is to be hoped that the theory will be improved still further on this important point. In \textsc{Bürg}'s lunar tables this equation was taken to be proportional to the sine of that argument diminished by three times the longitude of the Sun's perigee, with its coefficient assumed $= 14''$. For the horizontal parallax Mr. \textsc{Burckhardt} followed \textsc{Laplace}'s theory throughout. The ratio of that parallax to the Moon's horizontal semidiameter is based on the observed duration of meridian transits at full moon. --- A pair of new formulas developed by Mr. \textsc{Burckhardt} is appended at the end of the work. They can be used to advantage in calculating new and full moons.
